% Lecture notes: Week 13, Lecture 2 -- Thermal Expansion: Examples
% To hide this lecture from the website, add a line containing only:  % draft
% To set the date shown on the website, add a line like:  % posted: 2026-09-02
\documentclass[11pt]{article}
\usepackage{mech2030notes}
\begin{document}

% ##################################################################
%  WEEK 13, LECTURE 2
% ##################################################################
\lecture{13}{2}{Thermal Expansion: Examples}

\noindent
(These are the worked solutions to Problems~2 and~3 of the W13L1 in-class problems.)

\subsection*{Useful relations}
\begin{align*}
  \delta_T &= \alpha\,\Delta T\,L && \text{(thermal expansion, constant $\Delta T$)} \\
  \delta &= \frac{NL}{EA} && \text{(axial load, constant $N$, $E$, $A$)}
\end{align*}

% ==================================================================
\section*{Problem 2}
% ==================================================================

\noindent
\begin{minipage}[c]{0.55\textwidth}
\textbf{Setup:} The $\SI{50}{mm}$ diameter steel rod ($E = \SI{200}{GPa}$,
$\alpha = 12 \times 10^{-6}\ /{}^\circ\mathrm{C}$) is fixed at both ends. There is no load
in the members when $T_1 = \SI{80}{\celsius}$ and no load is applied.

\smallskip
\textbf{Find:} Determine the reactions at $A$ and $B$ when the load $P = \SI{200}{kN}$ is
applied and the temperature is reduced to $T_2 = \SI{20}{\celsius}$.
\end{minipage}%
\hfill
\begin{minipage}[c]{0.42\textwidth}
\centering
\begin{tikzpicture}
  \lwall{0}{-0.8}{0.8}
  \rwall{5}{-0.8}{0.8}
  \draw[beam] (0,-0.3) rectangle (5,0.3);
  \node[above left] at (-0.35,0.3) {$A$};
  \node[above right] at (5.35,0.3) {$B$};
  \node[pin] at (2.5,0) {}; \node[above] at (2.5,0.3) {$C$};
  \draw[force] (1.5,0) -- (2.45,0);
  \node[above] at (1.8,0.3) {$P$};
  \draw[dim] (0,-1.2) -- (2.5,-1.2) node[midway, fill=white] {$\SI{0.5}{m}$};
  \draw[dim] (2.5,-1.2) -- (5,-1.2) node[midway, fill=white] {$\SI{0.5}{m}$};
\end{tikzpicture}
\end{minipage}

\subsection*{FBD}

\begin{center}
\begin{tikzpicture}
  \draw[thick] (0,0) -- (5,0);
  \node[pin] at (0,0) {}; \node[pin] at (2.5,0) {}; \node[pin] at (5,0) {};
  \draw[force] (-1.2,0) node[left] {$A_x$} -- (-0.08,0);
  \draw[force] (2.5,-0.35) -- (3.5,-0.35) node[right] {$P = \SI{200}{kN}$};
  \draw[force] (5.08,0) -- (6.2,0) node[right] {$B_x$};
\end{tikzpicture}
\end{center}

\[
  \textstyle\sum F_x = 0 \quad\rightarrow\quad A_x + P + B_x = 0
  \qquad\Rightarrow\qquad
  \boxed{A_x = -P - B_x}
\]

\noindent
$\rightarrow$ Compatibility condition: \quad $\boxed{\delta_B = 0}$

\subsection*{Superposition}

\noindent
Remove the wall at $B$, and add up the displacement of $B$ due to each effect separately:

\begin{center}
\begin{tikzpicture}[scale=0.72, every node/.style={scale=0.9}]
  % original
  \begin{scope}
    \lwall{0}{-0.7}{0.7}
    \rwall{5}{-0.7}{0.7}
    \draw[beam] (0,-0.3) rectangle (5,0.3);
    \node[pin] at (2.5,0) {};
    \draw[force] (2.55,0) -- (3.5,0) node[right] {$P$};
  \end{scope}
  % (I) B_x
  \begin{scope}[shift={(0,-2.0)}]
    \node at (-1.0,0) {\Large$=$};
    \lwall{0}{-0.7}{0.7}
    \draw[beam] (0,-0.3) rectangle (5,0.3);
    \draw[thin, dashed] (5,0.3) -- (5.6,0.3) -- (5.6,-0.3) -- (5,-0.3);
    \node[pin] at (5,0) {};
    \draw[force] (5.05,0) -- (6.1,0) node[right] {$B_x$};
    \draw[{Bar[width=2mm]}-{Stealth[length=2mm]}, thin] (5,0.6) -- (5.6,0.6) node[right] {$\delta_B^{B_x}$};
    \draw[dim] (0,-0.75) -- (5,-0.75) node[midway, fill=white] {$L$};
    \node at (8.6,0) {$\left(\dfrac{B_x L}{EA}\right)$};
  \end{scope}
  % (II) thermal
  \begin{scope}[shift={(0,-4.2)}]
    \node at (-1.0,0) {\Large$+$};
    \lwall{0}{-0.7}{0.7}
    \draw[beam] (0,-0.3) rectangle (5,0.3);
    \draw[thin, dashed] (5,0.3) -- (5.6,0.3) -- (5.6,-0.3) -- (5,-0.3);
    \draw[{Bar[width=2mm]}-{Stealth[length=2mm]}, thin] (5,0.6) -- (5.6,0.6) node[right] {$\delta_B^{T}$};
    \node[below] at (2.5,-0.3) {\small$\Delta T = T_2 - T_1 = -\SI{60}{\celsius}$};
    \node at (8.6,0) {$\left(\alpha\,\Delta T\,L\right)$};
  \end{scope}
  % (III) P
  \begin{scope}[shift={(0,-6.4)}]
    \node at (-1.0,0) {\Large$+$};
    \lwall{0}{-0.7}{0.7}
    \draw[beam] (0,-0.3) rectangle (5,0.3);
    \draw[thin, dashed] (5,0.3) -- (5.6,0.3) -- (5.6,-0.3) -- (5,-0.3);
    \node[pin] at (2.5,0) {};
    \draw[force] (2.55,0) -- (3.5,0) node[right] {$P$};
    \draw[{Bar[width=2mm]}-{Stealth[length=2mm]}, thin] (5,0.6) -- (5.6,0.6) node[right] {$\delta_B^{P}$};
    \draw[dim] (0,-0.75) -- (2.5,-0.75) node[midway, fill=white] {$L/2$};
    \node at (8.6,0) {$\left(\dfrac{P\,(L/2)}{EA}\right)$};
  \end{scope}
\end{tikzpicture}
\end{center}

\noindent
$\rightarrow$ So,
\[
  \delta_B = \delta_B^{B_x} + \delta_B^{T} + \delta_B^{P}
  = \frac{B_x L}{EA} + \alpha\,\Delta T\,L + \frac{P\,(L/2)}{EA} = 0
\]
Dividing through by $L$:
\[
  \frac{B_x}{EA} = -\alpha\,\Delta T - \frac12\,\frac{P}{EA}
  \qquad\Rightarrow\qquad
  \boxed{B_x = -\alpha(\Delta T)\,EA - \tfrac12 P}
\]
\begin{align*}
  B_x &= -\left(12 \times 10^{-6}\ /{}^\circ\mathrm{C}\right)(-\SI{60}{\celsius})
         (\SI{200e9}{Pa})\left(\pi(\SI{0.025}{m})^2\right)
         - \tfrac12(\SI{200e3}{N})
\end{align*}
\[
  \fbox{\fbox{$B_x = \SI{182.7}{kN}$}}
\]
\[
  A_x = -P - B_x = -\SI{200}{kN} - \SI{182.7}{kN}
  \qquad\Rightarrow\qquad
  \fbox{\fbox{$A_x = -\SI{382.7}{kN}$}}
\]

\noindent
\textbf{Final FBD} (forces in their actual directions):

\begin{center}
\begin{tikzpicture}
  \draw[thick] (0,0) -- (5,0);
  \node[pin] at (0,0) {}; \node[pin] at (2.5,0) {}; \node[pin] at (5,0) {};
  \draw[force] (-0.08,0) -- (-1.2,0) node[left] {$\SI{382.7}{kN}$};
  \draw[force] (2.5,-0.35) -- (3.5,-0.35) node[right] {$\SI{200}{kN}$};
  \draw[force] (5.08,0) -- (6.2,0) node[right] {$\SI{182.7}{kN}$};
\end{tikzpicture}\\[0.3em]
(All segments are in \emph{tension}, due to $\Delta T < 0$.)
\end{center}

% ==================================================================
\section*{Problem 3}
% ==================================================================

\noindent
\textbf{Setup:} Use the same setup as Problem~2, but with $P$ unknown.

\smallskip
\noindent
\textbf{Find:} (a) What value of the applied load $P$ results in $B_x = 0$?
(b) What is $\delta_C$ in that case?

\subsection*{(a) Load for zero reaction at $B$}

\noindent
From Problem~2:
\[
  B_x = 0 = -\alpha(\Delta T)\,EA - \tfrac12 P
  \qquad\Rightarrow\qquad
  P = -2\,\alpha(\Delta T)\,EA
\]
\[
  \fbox{\fbox{$P = \SI{565}{kN}$}}
\]

\subsection*{(b) Displacement of $C$}

\noindent
$\rightarrow$ Superposition for $\delta_C$:

\begin{center}
\begin{tikzpicture}
  % thermal
  \begin{scope}
    \node at (-1.0,0) {$\delta_C \;=$};
    \lwall{0}{-0.7}{0.7}
    \draw[beam] (0,-0.3) rectangle (2.5,0.3);
    \draw[thin, dashed] (2.5,0.3) -- (5,0.3) -- (5,-0.3) -- (2.5,-0.3);
    \node[pin] at (2.5,0) {}; \node[left] at (2.45,0) {$C$};
    \draw[{Bar[width=2mm]}-{Stealth[length=2mm]}, thin] (2.5,-0.6) -- (3.1,-0.6) node[right] {$\delta_C^{T}$};
    \draw[dim] (0,0.65) -- (2.5,0.65) node[midway, fill=white] {$L/2$};
    \node[below] at (1.25,-0.3) {\small$\Delta T$};
    \node at (7.6,0) {$\left(\alpha(\Delta T)(L/2)\right)$};
  \end{scope}
  % P
  \begin{scope}[shift={(0,-2.2)}]
    \node at (-1.0,0) {\Large$+$};
    \lwall{0}{-0.7}{0.7}
    \draw[beam] (0,-0.3) rectangle (2.5,0.3);
    \draw[thin, dashed] (2.5,0.3) -- (5,0.3) -- (5,-0.3) -- (2.5,-0.3);
    \node[pin] at (2.5,0) {}; \node[left] at (2.45,0) {$C$};
    \draw[force] (2.55,0) -- (3.5,0) node[right] {$P$};
    \draw[{Bar[width=2mm]}-{Stealth[length=2mm]}, thin] (2.5,-0.6) -- (3.1,-0.6) node[right] {$\delta_C^{P}$};
    \draw[dim] (0,0.65) -- (2.5,0.65) node[midway, fill=white] {$L/2$};
    \node at (7.6,0) {$\left(\dfrac{P\,(L/2)}{EA}\right)$};
  \end{scope}
\end{tikzpicture}
\end{center}

\begin{itemize}[label=\then]
  \item No $\delta_C^{B_x}$ term, since $B_x = 0$.
\end{itemize}

\noindent
$\rightarrow$ So,
\begin{align*}
  \delta_C &= \alpha\,\Delta T\left(\frac{L}{2}\right) + \frac{P\,(L/2)}{EA} \\
  &= \left(12 \times 10^{-6}\ /{}^\circ\mathrm{C}\right)(-\SI{60}{\celsius})(\SI{0.5}{m})
   + \frac{(\SI{565e3}{N})(\SI{0.5}{m})}{(\SI{200e9}{Pa})\left(\pi(\SI{0.025}{m})^2\right)}
\end{align*}
\[
  \fbox{\fbox{$\delta_C = \SI{0.00036}{m} = \SI{0.36}{mm}$}} \quad \text{(rightwards)}
\]

\noindent
\textbf{Final FBD:}

\begin{center}
\begin{tikzpicture}
  \draw[thick] (0,0) -- (5,0);
  \node[pin] at (0,0) {}; \node[pin] at (2.5,0) {}; \node[pin] at (5,0) {};
  \node[above] at (0,0.05) {$A$};
  \node[above] at (2.5,0.05) {$C$};
  \node[above] at (5,0.05) {$B$};
  \draw[force] (-0.08,0) -- (-1.2,0) node[left] {$\SI{565}{kN}$};
  \draw[force] (2.5,-0.35) -- (3.5,-0.35) node[right] {$\SI{565}{kN}$};
\end{tikzpicture}\\[0.3em]
(No stress in segment $CB$!)
\end{center}

\end{document}
